% Seção 11: conclusões
A coupled formulation with displacement, pore-water and pore-air pressures was implemented in NeoPZ, validated against
the Liakopoulos drainage test and used to quantify how the hypotheses of the usual chain of analysis, uncoupled
seepage with atmospheric pore air followed by a stability analysis with the frozen pressures, bias the rainfall that
brings an unsaturated slope to failure. The main conclusions are the following.
\begin{enumerate}
\item With the failure criterion FS $\le1$, models must be compared by the critical intensity and the time to
failure, not by the minimum FS. A minimum FS below one is the FS of the slope saturated to its surface, a property
of the boundary-value problem that can be computed directly with the surface ponded.
\item For the homogeneous slope of \citet{rahardjo2007} with $k_s=36\mmh$ the usual model gives a critical intensity
of 48~mm/h ($1.3k_s$) for 24-h rainfalls, and the complete model 31~mm/h ($0.9k_s$) with an associated flow rule and
26~mm/h ($0.7k_s$) without dilatancy. The usual hypotheses overestimate the failure-triggering rainfall by 54--86\%
and delay the predicted failure by up to 10~h in intense rainfalls. The bias persists with $\chi=S_e$ (48 against
30~mm/h) and with a finer mesh (about 46 against 30~mm/h).
\item The two simplifications act in opposite directions. The trapped air halves the critical intensity; the coupling
raises it, by 4--10\% through the elastic swelling of the skeleton and by up to 63\% when the plastic dilation of an
associated flow rule is added. The air dominates, and a two-phase flow analysis on a rigid skeleton recovers the bias
with a critical intensity 4--20\% below that of the complete model.
\item The trapped air lowers the FS with less water. Once the surface ponds, the air between the wetting front and
the water table is compressed to 36--38~kPa; less water infiltrates and the matric suction increases, but the
pore-water pressure rises with the air pressure and the effective stress falls. A low-permeability slope
($k_s=3.6\mmh$, $63.4^\circ$) that the usual model places 9--10\% above failure reaches limit equilibrium; the effect
does not depend on whether the base and lateral boundaries are open to the air.
\item Dilatancy does not decide whether a critical intensity exists but shifts it by about 20\% and decides the
response after it is exceeded: without dilatancy the coupled transient loses equilibrium, independently of the time
step; with $\psi\ge\phi'/2$ the slope stays at limit equilibrium.
\item The bias concentrates in the regime $\Ic\gtrsim k_s$, in which the surface ponds and seals the air. For a
permeable soil whose critical intensity is well below $k_s$ the usual model is adequate (135 against 139~mm/h).
\end{enumerate}
Rainfall thresholds derived from uncoupled, passive-air analyses are therefore unconservative for slopes whose
critical rainfall exceeds the saturated conductivity of the surface soil. The validation of the air phase in
infiltration with measured air pressures, natural (open, heterogeneous) slopes and three-dimensional effects are the
next steps.
